% !TEX root = ../main.tex
% ============================================================================
% 5. AVALIAÇÃO CRÍTICA
% ============================================================================

\section{Avaliação Crítica da Estratégia de Testes}

A estratégia de testes do Educado é \textbf{insuficiente} diante do escopo do
produto analisado. A avaliação não decorre da ausência total de testes: há uma
base automatizada relevante no \texttt{educado-api}, um teste de fumaça no
\texttt{educado-web}, um teste de fumaça no \texttt{educado-app} e um ciclo de
avaliação de campo conduzido com usuários finais \cite{educado-repos,psp5}. O
problema é que essas frentes não formam uma proteção equilibrada para o
funcionamento do produto. O Web, analisado em maior profundidade neste
relatório, ainda tem cobertura automatizada mínima, e a integração entre painel,
API e experiência final do usuário não é verificada por testes de sistema ou
end-to-end \cite{educado-repos}.

Na execução realizada pela equipe, o backend apresentou 31 suítes e 506 testes
aprovados. O relatório de cobertura do Jest indicou 79,36\% de instruções,
69,26\% de branches, 78\% de funções e 78,9\% de linhas no escopo configurado
em \texttt{jest.config.ts}; esses números dependem da configuração de cobertura
do próprio repositório e não representam todo o produto \cite{educado-repos}. No
\texttt{educado-web} e no \texttt{educado-app}, foi executado somente um teste
automatizado em cada repositório \cite{educado-repos}. A frente de campo, por
outro lado, mede usabilidade e jornada real com catadores em dispositivo Android,
mas não se conecta à suíte automatizada nem gera uma barreira de regressão no CI
\cite{psp5,dashboard}.

Assim, o veredito é insuficiente: o projeto possui peças úteis, mas a estratégia
como conjunto ainda não garante que mudanças no Web, na API ou na experiência
final preservem fluxos críticos, metas de qualidade e comportamentos observados
em campo.


% ----------------------------------------------------------------------------
% 5.1 PONTOS FORTES
% ----------------------------------------------------------------------------

\subsection{Pontos Fortes}

Os pontos fortes mostram que o projeto já tem material para evoluir: há
automação no repositório, documentação de contribuição e uma frente de validação
com usuários. Eles não bastam para tornar a estratégia adequada, mas reduzem o
custo de amadurecê-la.

{\footnotesize
\begin{longtable}{>{\raggedright\arraybackslash}p{2.9cm} >{\raggedright\arraybackslash}p{5.1cm} >{\raggedright\arraybackslash}p{5.1cm}}
\caption{Pontos fortes da estratégia de testes.}
\label{tab:pontosfortes}\\
\toprule
\textbf{Aspecto} & \textbf{Evidência} & \textbf{Justificativa} \\
\midrule
\endfirsthead

\toprule
\textbf{Aspecto} & \textbf{Evidência} & \textbf{Justificativa} \\
\midrule
\endhead

Suíte unitária relevante na API
&
A execução de \texttt{npm test} no \texttt{educado-api} resultou em 31 suítes e
506 testes aprovados. Os testes estão distribuídos por módulos de aplicação,
segurança, armazenamento e middlewares \cite{educado-repos}.
&
A quantidade de casos e sua distribuição por diferentes domínios fornecem uma
base real de proteção para regras de negócio, validações, autenticação,
autorização, recuperação de senha, mídia, cursos, matrículas e progresso dos
estudantes.
\\
\addlinespace

Cobertura mensurável na API
&
O script \texttt{test:coverage} da API produz métricas de cobertura e o
\texttt{jest.config.ts} delimita os arquivos considerados na medição
\cite{educado-repos}.
&
A coleta de cobertura oferece uma régua objetiva para acompanhar evolução da
suíte. A ressalva é importante: por ser delimitada pela configuração, a métrica
é útil para a API, mas não deve ser lida como cobertura global do produto.
\\
\addlinespace

Execução automática em pull requests
&
Os workflows dos repositórios \texttt{educado-api}, \texttt{educado-web} e
\texttt{educado-app} executam testes em pull requests, junto com verificações de
tipo e, conforme o repositório, lint ou build \cite{educado-repos}.
&
A automação reduz a dependência de execução manual e cria um ponto natural para
detectar regressões antes da integração de uma mudança. O problema remanescente
é que esse aviso ainda não funciona como portão obrigatório de qualidade.
\\
\addlinespace

Organização por domínio na API
&
No \texttt{educado-api}, os testes ficam próximos ao código testado, em
diretórios \texttt{\_\_tests\_\_}, e seguem a nomenclatura \texttt{*.test.ts}
\cite{educado-repos}.
&
A organização facilita localização e manutenção. O uso de mocks também torna a
suíte rápida e determinística; esse é o lado positivo do mesmo tradeoff que, nas
limitações, aparece como risco de divergência em relação às dependências reais.
\\
\addlinespace

Processo de contribuição orientado à qualidade
&
O guia de contribuição e os modelos de pull request pedem descrição de testes,
execução de verificações e inclusão ou atualização de testes para mudanças de
comportamento \cite{educado-repos}.
&
Esses critérios incentivam quem contribui a explicitar como verificou a mudança
e ajudam a revisão humana. Como a própria documentação exige fonte e evidência,
ela também cria um padrão útil para avaliar as lacunas encontradas no relatório.
\\
\addlinespace

Avaliação de campo com usuários reais
&
O ciclo de campo avaliou a experiência do Educado com catadores, usando roteiro
de usabilidade, indicadores e segunda visita para verificar melhorias
\cite{psp5,dashboard}.
&
Essa frente cobre aspectos que a suíte automatizada não mede, como compreensão
da interface, autonomia de uso, desempenho percebido e aderência da jornada ao
contexto real dos usuários.
\\
\addlinespace
\bottomrule
\end{longtable}
}


% ----------------------------------------------------------------------------
% 5.2 LIMITAÇÕES
% ----------------------------------------------------------------------------

\subsection{Limitações e Problemas}

As limitações tornam a estratégia insuficiente porque atingem justamente o
comportamento integrado do produto. A API tem a parte mais madura da automação,
mas o Web, que é o foco principal do recorte, permanece quase sem testes de
comportamento, e a frente de campo não se transforma em critério automatizado de
regressão.

{\footnotesize
\begin{longtable}{>{\raggedright\arraybackslash}p{2.9cm} >{\raggedright\arraybackslash}p{5.1cm} >{\raggedright\arraybackslash}p{5.1cm}}
\caption{Limitações identificadas.}
\label{tab:limitacoes}\\
\toprule
\textbf{Limitação} & \textbf{Evidência} & \textbf{Impacto} \\
\midrule
\endfirsthead

\toprule
\textbf{Limitação} & \textbf{Evidência} & \textbf{Impacto} \\
\midrule
\endhead

Baixa proteção automatizada do Web
&
O \texttt{educado-web} contém apenas o teste de fumaça
do seletor de idioma; não há suíte cobrindo autenticação, permissões,
formulários, criação de cursos ou integração com a API \cite{educado-repos}.
&
Funcionalidades administrativas e de criação de conteúdo podem regredir sem que
o CI detecte. Como o relatório analisa o Web em maior profundidade, essa lacuna
pesa mais que a quantidade de testes existentes na API.
\\
\addlinespace

Integração e end-to-end ausentes
&
Não foram encontradas configurações ou suítes com Playwright, Cypress, Detox ou
Maestro. O \texttt{educado-api} possui Supertest entre as dependências, mas não
há testes HTTP usando a biblioteca \cite{educado-repos}.
&
Os testes existentes não validam fluxos completos entre interface, API,
autenticação, persistência e serviços externos. Um módulo pode funcionar
isoladamente e falhar quando integrado ao restante do sistema.
\\
\addlinespace

Predomínio de mocks na API
&
Diversos testes substituem modelos, JWT, armazenamento S3 e outros
colaboradores por \texttt{jest.mock} e \texttt{jest.fn}; não há suíte equivalente
com PostgreSQL, Redis e MinIO reais ou efêmeros \cite{educado-repos}.
&
Mocks tornam a suíte rápida, mas podem divergir do comportamento real. O texto
trata isso como tradeoff: é força para feedback curto e limitação para confiança
em consultas, transações, serialização, infraestrutura e contratos externos.
\\
\addlinespace

Cobertura sem critério de aprovação
&
O workflow da API executa \texttt{npm test}, mas não executa
\texttt{npm run test:coverage}. O \texttt{jest.config.ts} não define
\texttt{coverageThreshold}; Web e App também não possuem metas versionadas
\cite{educado-repos}.
&
A cobertura pode cair sem quebrar o pipeline. O crescimento da suíte depende de
revisão manual, sem limite objetivo para impedir integração de código crítico
sem testes suficientes.
\\
\addlinespace

Documentação inconsistente
&
Documentos do Web e do App afirmam não haver suíte automatizada, embora ambos
tenham teste e workflow executando \texttt{npm test}. A documentação da API
também descreve testes de integração que não aparecem no código
\cite{educado-repos}.
&
A divergência reduz a confiança na documentação e pode induzir contribuidores a
ignorar infraestrutura existente ou acreditar em uma proteção que ainda não foi
implementada.
\\
\addlinespace

Frente de campo desconectada do CI
&
O PSP5 e o dashboard registram indicadores, melhorias e segunda visita de
validação, mas esses achados não aparecem como testes automatizados, critérios
de aceite versionados ou gates de regressão nos repositórios
\cite{psp5,dashboard,educado-repos}.
&
Problemas reais de usabilidade, clareza, desempenho percebido e autonomia podem
reaparecer sem alerta automático. A estratégia mede o usuário em campo, mas não
transforma essa aprendizagem em proteção contínua do produto.
\\
\addlinespace

Execução restrita a pull requests
&
Os workflows de teste dos três repositórios usam o evento
\texttt{pull\_request}; não foram identificadas execuções equivalentes em pushes
para a branch principal, tags de versão ou rotinas agendadas
\cite{educado-repos}.
&
Falhas introduzidas por merges, mudanças de dependências ou diferenças no
processo de release podem não ser verificadas novamente depois da aprovação do
pull request.
\\
\addlinespace
\bottomrule
\end{longtable}
}


% ----------------------------------------------------------------------------
% 5.3 RECOMENDAÇÕES
% ----------------------------------------------------------------------------

\subsection{Recomendações de Melhoria}

As recomendações priorizam o Web e a integração com a API, mantendo o App e a
frente de campo como contexto necessário para a estratégia completa. Não se
recomenda buscar apenas um número alto de testes, mas selecionar casos por risco,
frequência de uso e impacto para estudantes, criadores de conteúdo,
administradores e usuários finais.

{\footnotesize
\begin{longtable}{>{\raggedright\arraybackslash}p{2.4cm} >{\raggedright\arraybackslash}p{3.8cm} >{\raggedright\arraybackslash}p{4.5cm} >{\raggedright\arraybackslash}p{1.9cm}}
\caption{Recomendações de melhoria.}
\label{tab:recomendacoes}\\
\toprule
\textbf{Problema} & \textbf{Recomendação} & \textbf{Justificativa} & \textbf{Prioridade} \\
\midrule
\endfirsthead

\toprule
\textbf{Problema} & \textbf{Recomendação} & \textbf{Justificativa} & \textbf{Prioridade} \\
\midrule
\endhead

Cobertura mínima no Web
&
Criar testes unitários e de componente com Vitest, jsdom e Testing Library para
autenticação, roteamento, internacionalização, formulários, permissões, cliente
HTTP e criação ou edição de cursos.
&
A infraestrutura do Web já existe, mas só há teste de fumaça
\cite{educado-repos}. Ampliar essa suíte ataca a lacuna mais importante dentro
do recorte principal do relatório.
&
Alta
\\
\addlinespace

Integração real com a API
&
Adicionar testes HTTP com Supertest na API e cenários de integração do Web contra
contratos estáveis da API. Quando necessário, usar infraestrutura isolada para
PostgreSQL, Redis e MinIO.
&
A API concentra a suíte mais madura, mas ainda depende fortemente de mocks e não
exercita rotas HTTP com Supertest \cite{educado-repos}. Testes de integração
reduzem o risco de falhas entre painel, autenticação, persistência e serviços.
&
Alta
\\
\addlinespace

Ausência de end-to-end
&
Implementar uma suíte pequena de end-to-end para os fluxos Web mais críticos,
preferencialmente com Playwright: login, acesso por perfil, criação de curso,
publicação e consumo básico refletido na API.
&
Poucos fluxos bem escolhidos validam o produto como sistema sem transformar o CI
em uma suíte lenta e frágil. Para o App, ferramentas como Maestro ou Detox devem
ser tratadas como evolução posterior, pois o foco deste relatório é o Web.
&
Alta
\\
\addlinespace

Achados de campo sem regressão
&
Transformar indicadores e problemas recorrentes do PSP5 em critérios de aceite,
checklists versionados e, quando possível, testes automatizados de interface.
&
A avaliação com catadores identificou metas e problemas reais de uso
\cite{psp5,dashboard}. Incorporar esses achados impede que a frente de campo
fique isolada da evolução técnica do produto.
&
Alta
\\
\addlinespace

Cobertura sem limite mínimo
&
Executar cobertura no CI e estabelecer metas progressivas, começando próximas
aos valores atuais da API. Definir também metas para código novo ou modificado e
ampliar gradualmente o \texttt{collectCoverageFrom}.
&
Um limite progressivo evita queda de cobertura sem bloquear imediatamente o
projeto por módulos legados. A análise de código novo torna o critério mais
viável e mais justo.
&
Média
\\
\addlinespace

Documentação divergente
&
Atualizar README, CONTRIBUTING e modelos de pull request para refletir os testes
existentes, os comandos atuais e a exigência de fonte para afirmações técnicas.
&
A documentação já orienta contribuições, mas está desalinhada em pontos
importantes \cite{educado-repos}. Corrigir isso reduz erro de contribuição e
ajuda revisores a cobrar evidências de forma consistente.
&
Média
\\
\addlinespace

CI limitado a pull requests
&
Executar a suíte também em pushes para branches principais, além de incluir
verificação em releases. Testes mais caros podem ficar em jobs separados ou
agendados.
&
Executar só em pull request deixa lacunas após merge e durante release
\cite{educado-repos}. Separar suítes por custo preserva feedback rápido sem
abrir mão de verificação periódica mais profunda.
&
Média
\\
\addlinespace
\bottomrule
\end{longtable}
}

Em síntese, a estratégia atual tem bons pontos de partida, mas ainda é
insuficiente. A prioridade deve ser proteger o Web com testes de comportamento,
validar sua integração com a API, transformar os achados de campo em critérios de
regressão e fazer com que o CI deixe de apenas executar verificações para
sustentar uma política objetiva de qualidade.
